\scanpage{46}
\mspage{2}
% archive sheet number: 45
\uline{objets}: Triples $(X,\widetilde{\Sigma},\deleted{\unsure{S|\delta}}R)$, avec
\margin{$X,\deleted{\unsure{\widetilde{\Sigma}}}\partial X,\Sigma,S$}
\[
\begin{aligned}
& X\ \text{surface à bord compacte orientée}\\
& \widetilde{\Sigma}\ \text{sous-variété compacte de }X\ \deleted{\unsure{\partial X|XXX}},\quad \widetilde{\Sigma}\mathrel{\unsure{\supset|>}}\partial X\\
& \deleted{\unsure{S|\delta}}R\subset\Sigma\ \deleted{\unsure{\setminus\partial X}},\ \text{partie finie}
\end{aligned}
\]
% slanted
\margin{$\widetilde{\Sigma}=\partial X\amalg\Sigma$\\On peut considérer que $R$ rigidifie, mais \unsure{non} que $R$ ``\unsure{découpe}'' $\Sigma\ldots$}
\uline{flèches}: classes d'isotopie d'homéom. conservant l'orientation

\deleted{Variante \unsure{avec}} Les composantes connexes de $\widetilde{\Sigma}\setminus\partial X$ \inserted{$=\Sigma$} \unsure{sont} les \uline{cercles de découpage}, \inserted{celles de $\partial X$ \unsure{sont} les \uline{cercles de raccord},} les éléments de $\deleted{\unsure{S|\delta}}R$ \unsure{sont} les \deleted{\unsure{points marqués|points repères}} \uline{sommets de \deleted{\unsure{marquage|repérage}}\inserted{rigidification}}, \deleted{\unsure{(ou de rigidification)}}, \deleted{les composantes connexes de $\partial X$ sont les ``\uline{\unsure{trous}}''}. Les composantes connexes de $X$ sont les \deleted{\unsure{morceaux}}\inserted{\uline{\unsure{composants|composantes}}} de l'objet envisagé.

Variante
\[
S_{*}\unsure{\mathscr{T}|\mathcal{T}}D_{\mathrm{top}}^{\pm},\quad \text{avec}
\]
\diagram{
\begin{tikzcd}[ampersand replacement=\&]
S_{*}\mathscr{T}D_{\mathrm{top}}\arrow[r,hook,no head] \& S_{*}\mathscr{T}D_{\mathrm{top}}^{\pm}\arrow[r] \& [\pm1]
\end{tikzcd}
}
comme tantôt.

On a des sous-groupoïdes pleins
\diagram{
\(\begin{array}{cl}
\begin{gathered}
R=\unsure{\varnothing|\phi},\quad \unsure{\Sigma|\bar{\Sigma}}=\deleted{\unsure{\partial X}}\varnothing\\
\deleted{\unsure{\Sigma=\varnothing}}
\end{gathered}
 &\mathscr{T}_{\mathrm{top}}\ \text{et}\ \mathscr{T}_{\mathrm{top}}^{\pm}\\
\begin{gathered}
\Sigma=\deleted{\unsure{\varnothing}}\partial X\\
R\ \text{quelconque}
\end{gathered}
 & S_{*}\mathscr{T}_{\mathrm{top}}\ \text{et}\ S_{*}\mathscr{T}_{\mathrm{top}}^{\pm}\\
\deleted{\unsure{S}}R=\varnothing & \mathscr{T}\deleted{\unsure{{}_{*}}}D_{\mathrm{top}}\ \text{et}\ \mathscr{T}D_{\mathrm{top}}^{\pm}
\end{array}\)
}
(groupoïdes de Teichmüller \uline{spéciaux} \inserted{(ou \uline{rigidifiés})} (\unsure{topologiques|topologique} orienté))

(Cas pas \unsure{intéressant} je crois, sauf \deleted{pour} $\Sigma\deleted{\unsure{=}}\unsure{\ne|=}\varnothing$ également)

% slanted
\margin{$MD=$ ``modulaire'' ou Mumford-Deligne}
\deleted{\unsure{Dé}} \uline{Détermination} du $\pi_{0}(S_{*}\mathscr{T}D_{\mathrm{top}})\xrightarrow{\sim}\pi_{0}(S_{*}\mathscr{T}D_{\mathrm{top}}^{\pm})$.

On introduit les \uline{graphes} \deleted{\unsure{MD}} bipondérés
